\documentclass[11pt,a4paper]{article}

\usepackage[margin=0.72in]{geometry}
\usepackage{libertinus}
\usepackage{microtype}
\usepackage{xcolor}
\usepackage{enumitem}
\usepackage{tabularx}
\usepackage{array}
\usepackage{booktabs}
\usepackage{longtable}
\usepackage{amsmath}
\usepackage{amssymb}
\usepackage{fancyhdr}
\usepackage{titlesec}
\usepackage{tcolorbox}
\usepackage{hyperref}
\usepackage{multicol}
\usepackage{parskip}

\definecolor{navy}{HTML}{183153}
\definecolor{blue}{HTML}{2563EB}
\definecolor{green}{HTML}{166534}
\definecolor{red}{HTML}{B91C1C}
\definecolor{gold}{HTML}{92400E}
\definecolor{lightblue}{HTML}{EFF6FF}
\definecolor{lightgreen}{HTML}{F0FDF4}
\definecolor{lightred}{HTML}{FEF2F2}
\definecolor{lightgold}{HTML}{FFFBEB}

\hypersetup{colorlinks=true,linkcolor=blue,urlcolor=blue}
\pagestyle{fancy}
\fancyhf{}
\lhead{SUBJECT -- MODULE / UNIT}
\rhead{Exam Study Notes}
\cfoot{\thepage}

\titleformat{\section}{\Large\bfseries\color{navy}}{}{0pt}{}
\titleformat{\subsection}{\large\bfseries\color{navy}}{}{0pt}{}
\titleformat{\subsubsection}{
\normalsize\bfseries\color{blue}}{}{0pt}{}

\setlist[itemize]{leftmargin=*,nosep}
\setlist[enumerate]{leftmargin=*,nosep}


\newtcolorbox{exambox}{colback=lightblue,colframe=blue,title=Exam Focus,fonttitle=\bfseries}

\newtcolorbox{rememberbox}{colback=lightgold,colframe=gold,title=Remember,fonttitle=\bfseries}

\newtcolorbox{dangerbox}{colback=lightred,colframe=red,title=Do Not Confuse,fonttitle=\bfseries}

\newtcolorbox{answerbox}{colback=lightgreen,colframe=green,title=Answer Skeleton,fonttitle=\bfseries}


\newcommand{\mkmarks}[1]{\hfill\textbf{[#1 marks]}}

\newcommand{\cmd}[1]{\texttt{#1}}

\begin{document}

\begin{center}
{\Huge\bfseries\color{navy} SUBJECT TITLE}\\[4pt]
{\LARGE\bfseries COURSE / MODULE / UNIT}\\[8pt]
{\large High-yield notes, comparisons, diagrams, and question bank}\\[10pt]
\rule{0.85\textwidth}{0.6pt}
\end{center}

\begin{exambox}
\textbf{Exam strategy.} [PLACEHOLDER: Add the exam pattern, marks distribution, priority guidance, and study strategy here.]
\end{exambox}

\tableofcontents

\newpage

\section{1. MAIN TOPIC}

\subsection{1.1 Subtopic}
[PLACEHOLDER: Definition / core concept / explanation.]

\begin{center}
\fbox{PLACEHOLDER DIAGRAM / FLOW / ARCHITECTURE}
\end{center}

\begin{rememberbox}
\textbf{Remember:} [PLACEHOLDER: one high-yield memory point.]
\end{rememberbox}

\subsection{1.2 Another Subtopic}
\begin{enumerate}
    \item [PLACEHOLDER: Step / point 1]
    \item [PLACEHOLDER: Step / point 2]
    \item [PLACEHOLDER: Step / point 3]
\end{enumerate}

\begin{dangerbox}
\textbf{Do not confuse:} [PLACEHOLDER: commonly confused terms or trap.]
\end{dangerbox}

\begin{exambox}
\textbf{Likely question:} [PLACEHOLDER: exam question and marks.]
\end{exambox}

\section{2. COMPARISON / CLASSIFICATION}

\subsection{2.1 Comparison Table}
\begin{center}
\small
\begin{tabularx}{\textwidth}{>{\bfseries}p{0.22\textwidth} X X}
\toprule
Feature & Concept A & Concept B \\ 
\midrule
Definition & [PLACEHOLDER] & [PLACEHOLDER] \\ 
Working & [PLACEHOLDER] & [PLACEHOLDER] \\ 
Advantages & [PLACEHOLDER] & [PLACEHOLDER] \\ 
Limitations & [PLACEHOLDER] & [PLACEHOLDER] \\ 
\bottomrule
\end{tabularx}
\end{center}

\section{3. MECHANISM / ALGORITHM / PROCESS}

\subsection{3.1 Working}
\begin{enumerate}
    \item [PLACEHOLDER: Step 1]
    \item [PLACEHOLDER: Step 2]
    \item [PLACEHOLDER: Step 3]
    \item [PLACEHOLDER: Step 4]
\end{enumerate}

\begin{center}
\texttt{PLACEHOLDER: process / timeline / packet / architecture diagram}
\end{center}

\subsection{3.2 Formula / Technical Detail}
\[
    \text{PLACEHOLDER FORMULA}
\]

\section{4. IMPORTANT EXAM TOPICS}

\begin{itemize}
    \item [PLACEHOLDER: High-yield topic]
    \item [PLACEHOLDER: High-yield topic]
    \item [PLACEHOLDER: High-yield topic]
\end{itemize}

\section{5. 10-Mark Answer Templates}

\subsection{Question A: [PLACEHOLDER QUESTION] \mkmarks{10}}
\begin{answerbox}
Write in this order:
\begin{enumerate}
    \item [PLACEHOLDER: Introduction / definition]
    \item [PLACEHOLDER: Main concept 1]
    \item [PLACEHOLDER: Main concept 2]
    \item [PLACEHOLDER: Mechanism / diagram]
    \item [PLACEHOLDER: Comparison / example]
    \item [PLACEHOLDER: Conclusion]
\end{enumerate}
\end{answerbox}

\subsection{Question B: [PLACEHOLDER QUESTION] \mkmarks{10}}
\begin{answerbox}
[PLACEHOLDER: Answer structure / points to include / diagram to draw.]
\end{answerbox}

\section{6. Question Bank}

\subsection*{Definitions / 2--3 marks}
\begin{enumerate}
    \item [PLACEHOLDER QUESTION]
    \item [PLACEHOLDER QUESTION]
    \item [PLACEHOLDER QUESTION]
\end{enumerate}

\subsection*{5-mark questions}
\begin{enumerate}
    \item [PLACEHOLDER QUESTION]
    \item [PLACEHOLDER QUESTION]
    \item [PLACEHOLDER QUESTION]
\end{enumerate}

\subsection*{10-mark questions}
\begin{enumerate}
    \item [PLACEHOLDER QUESTION]
    \item [PLACEHOLDER QUESTION]
    \item [PLACEHOLDER QUESTION]
\end{enumerate}

\section{7. Last-Minute Revision Page}

\begin{multicols}{2}
\textbf{CORE CONCEPT}\\
\textit{PLACEHOLDER: one-line summary.}

\textbf{KEY TERMS}\\
\textit{PLACEHOLDER: compact definitions.}

\textbf{PROCESS}\\
\textit{PLACEHOLDER: compressed mechanism.}

\textbf{FORMULAS}\\
\textit{PLACEHOLDER: formulas / relations.}

\textbf{DO NOT CONFUSE}\\
\textit{PLACEHOLDER: common traps.}

\textbf{EXAM DIAGRAMS}\\
\textit{PLACEHOLDER: diagrams to memorize/draw.}
\end{multicols}

\section*{Source Scope}
\textit{PLACEHOLDER: List the supplied source files and state that terminology, organization, and emphasis should follow the source material.}

\vfill
\begin{center}
\textit{Study principle: understand the mechanism first, then memorize the exam wording.}
\end{center}

\end{document}
